\begin{lemma}
\label{lemma-argument-proves}
Let $X$ be a quasi-compact and quasi-separated scheme. Suppose that
for every affine open $U \subset X$ the right derived functor
$$
\Phi : D(\QCoh(\mathcal{O}_U)) \to D(\QCoh(\mathcal{O}_X))
$$
of the left exact functor
$j_* : \QCoh(\mathcal{O}_U) \to \QCoh(\mathcal{O}_X)$
fits into a commutative diagram
$$
\xymatrix{
D(\QCoh(\mathcal{O}_U)) \ar[d]_\Phi \ar[r]_{i_U} &
D_\QCoh(\mathcal{O}_U) \ar[d]^{Rj_*} \\
D(\QCoh(\mathcal{O}_X)) \ar[r]^{i_X} &
D_\QCoh(\mathcal{O}_X)
}
$$
Then the functor (\ref{equation-compare})
$$
D(\QCoh(\mathcal{O}_X))
\longrightarrow
D_\QCoh(\mathcal{O}_X)
$$
is an equivalence with quasi-inverse given by $RQ_X$.
\end{lemma}

\begin{proof}
Let $E$ be an object of $D_\QCoh(\mathcal{O}_X)$ and
let $A$ be an object of $D(\QCoh(\mathcal{O}_X))$.
We have to show that the adjunction maps
$$
A \to RQ_X(i_X(A))
\quad\text{and}\quad
i_X(RQ_X(E)) \to E
$$
are isomorphisms. Consider the hypothesis
$H_n$: the adjunction maps above are isomorphisms
whenever $E$ and $i_X(A)$ are supported
(Definition \ref{definition-supported-on})
on a closed subset of $X$ which
is contained in the union of $n$ affine opens of $X$.
We will prove $H_n$ by induction on $n$.

\medskip\noindent
Base case: $n = 0$. In this case $E = 0$, hence the map
$i_X(RQ_X(E)) \to E$ is an isomorphism. Similarly $i_X(A) = 0$.
Thus the cohomology sheaves of $i_X(A)$ are zero. Since the inclusion
functor $\QCoh(\mathcal{O}_X) \to \textit{Mod}(\mathcal{O}_X)$
is fully faithful and exact, we conclude that the cohomology
objects of $A$ are zero, i.e., $A = 0$ and
$A \to RQ_X(i_X(A))$ is an isomorphism as well.

\medskip\noindent
Induction step. Suppose that $E$ and $i_X(A)$ are supported on a
closed subset $T$ of $X$ contained in $U_1 \cup \ldots \cup U_n$
with $U_i \subset X$ affine open. Set $U = U_n$.
Consider the distinguished triangles
$$
A \to \Phi(A|_U) \to A' \to A[1]
\quad\text{and}\quad
E \to Rj_*(E|_U) \to E' \to E[1]
$$
where $\Phi$ is as in the statement of the lemma.
Note that $E \to Rj_*(E|_U)$ is a quasi-isomorphism over $U = U_n$.
Since $i_X \circ \Phi = Rj_* \circ i_U$ by assumption
and since $i_X(A)|_U = i_U(A|_U)$
we see that $i_X(A) \to i_X(\Phi(A|_U))$ is a quasi-isomorphism over $U$.
Hence $i_X(A')$ and $E'$ are supported on the closed
subset $T \setminus U$ of $X$ which is contained in
$U_1 \cup \ldots \cup U_{n - 1}$.
By induction hypothesis the statement is true for $A'$ and $E'$. By
Derived Categories, Lemma \ref{derived-lemma-third-isomorphism-triangle}
it suffices to prove the maps
$$
\Phi(A|_U) \to RQ_X(i_X(\Phi(A|_U)))
\quad\text{and}\quad
i_X(RQ_X(Rj_*E|_U)) \to Rj_*(E|_U)
$$
are isomorphisms. By assumption and by
Lemma \ref{lemma-flat-pushforward-coherator}
(the inclusion morphism $j : U \to X$ is flat, quasi-compact, and
quasi-separated) we have
$$
RQ_X(i_X(\Phi(A|_U))) = RQ_X(Rj_*(i_U(A|_U))) = \Phi(RQ_U(i_U(A|_U)))
$$
and
$$
i_X(RQ_X(Rj_*(E|_U))) = i_X(\Phi(RQ_U(E|_U))) = Rj_*(i_U(RQ_U(E|_U)))
$$
Finally, the maps
$$
A|_U \to RQ_U(i_U(A|_U))
\quad\text{and}\quad
i_U(RQ_U(E|_U)) \to E|_U
$$
are isomorphisms by Lemma \ref{lemma-affine-coherator}. The result follows.
\end{proof}
